\documentclass[11pt]{article}

\usepackage[margin=0.7in]{geometry}
\usepackage[T1]{fontenc}
\usepackage{lmodern}
\usepackage{microtype}
\usepackage{amsmath,amssymb}
\usepackage{xcolor}
\usepackage{enumitem}
\usepackage[most]{tcolorbox}
\usepackage[hidelinks]{hyperref}

\hypersetup{
  pdftitle={Homework 5: Ring Theory, MATH 327, Fall 2026},
  pdfauthor={Manish Patnaik}
}

\definecolor{courseblue}{HTML}{1F4E79}
\definecolor{lightblue}{HTML}{EAF2F8}
\definecolor{boxborder}{HTML}{B8CBDC}
\definecolor{darkgray}{HTML}{333333}

\setlength{\parindent}{0pt}
\setlength{\parskip}{0.5em}
\setlist[enumerate]{leftmargin=2em, itemsep=0.25em, topsep=0.3em}

\newcommand{\Z}{\mathbb{Z}}
\newcommand{\F}{\mathbb{F}}

\newtcolorbox{exercise}[1]{
  enhanced,
  colback=white,
  colframe=boxborder,
  colbacktitle=lightblue,
  coltitle=courseblue,
  fonttitle=\bfseries,
  title={#1},
  boxrule=0.6pt,
  arc=1mm,
  left=2.5mm,
  right=2.5mm,
  top=1.2mm,
  bottom=1.2mm,
  before skip=0.55em,
  after skip=0.55em
}

\begin{document}

\begin{center}
  \colorbox{lightblue}{%
    \begin{minipage}{0.92\textwidth}
      \vspace{0.6em}
      \begin{tabular*}{\textwidth}{@{\extracolsep{\fill}}lr}
        {\color{courseblue}\bfseries\LARGE Homework 5} &
        {\color{courseblue}\bfseries\LARGE MATH 327}\\[0.4em]
        \multicolumn{2}{l}{\color{darkgray}\bfseries Ring theory $\cdot$ Fall 2026}\\[0.35em]
        \multicolumn{2}{l}{\color{darkgray}Lectures 25--29 $\cdot$ Suggested completion: Monday, November 30}
      \end{tabular*}
      \vspace{0.5em}
    \end{minipage}}
\end{center}

Work through each exercise after the indicated lecture. This set is for
practice. References are to Chapter 11 of Michael Artin's \emph{Algebra},
second edition.

\begin{exercise}{1. Units in polynomial rings (after Lecture 26)}
\begin{enumerate}[label=\textbf{(\alph*)}]
  \item Let $F$ be a field. Prove that the units of $F[x]$ are exactly the
        nonzero constant polynomials.
  \item Find an inverse of $1+2x$ in $(\Z/4\Z)[x]$. Explain why the degree
        argument from part (a) does not apply in this ring.
\end{enumerate}
\end{exercise}

\begin{exercise}{2. Intersections and generated ideals (after Lecture 27)}
\begin{enumerate}[label=\textbf{(\alph*)}]
  \item Let $I$ and $J$ be ideals of a ring $R$. Show that $I\cap J$ is an ideal.
  \item In $\Z$, give a single generator for each of $(6)\cap(10)$ and
        $(6,10)$. Justify both answers.
\end{enumerate}
{\small\color{darkgray}Adapted from Exercise 3.13.}
\end{exercise}

\begin{exercise}{3. Evaluation and the first isomorphism theorem (after Lecture 28)}
\begin{enumerate}[label=\textbf{(\alph*)}]
  \item Let $R$ be a ring and let $a\in R$. Use evaluation at $a$ to show
        that $R[x]/(x-a)\cong R$.
  \item Find the kernel and image of $\Z[x]\to\Z/5\Z$ given by
        $f(x)\mapsto f(2)\bmod 5$. What quotient ring does this identify?
\end{enumerate}
{\small\color{darkgray}Part (a) is Exercise 5.2.}
\end{exercise}

\begin{exercise}{4. Quotienting in stages (after Lecture 29)}
Identify each quotient as a quotient of $\F_p[x]$ for a suitable prime $p$.
Simplify where possible and give the number of elements.
\[
  \textbf{(a)}\ \Z[x]/(x^2-3,\,2x+4)
  \qquad
  \textbf{(b)}\ \Z[x]/(4x-5,\,2x^2-4)
  \qquad
  \textbf{(c)}\ \Z[x]/(5,\,x^2+3).
\]
{\small\color{darkgray}Adapted from Exercise 4.3.}
\end{exercise}

\begin{exercise}{5. Distinguishing two rings (after Lecture 29)}
Show that $\Z[x]/(x^2+7)$ and $\Z[x]/(2x^2+7)$ are not isomorphic.
\textit{Hint.} What must a ring isomorphism do to $2$ and to the ideal it
generates?
\par
{\small\color{darkgray}Exercise 4.4.}
\end{exercise}

\end{document}
